\chapter{Peano 存在定理}

\begin{solution}{peano-euler-stays-rectangle}
在第 \(k\) 个节点前，各段斜率范数至多 \(M\)，所以
\[
 \|y_n(t_k)-y_0\|
 \le\sum_{j<k}M\Delta=Mk\Delta.
\]
在段内再加 \(M(t-t_k)\)，得
\(\|y_n(t)-y_0\|\le M(t-t_0)\le Mh\le b\)。由归纳，每次计算斜率时
\((t_k,y_n(t_k))\) 都仍在 \(Q\)。
\end{solution}

\begin{solution}{peano-euler-equilipschitz}
若 \(s<t\)，把 \([s,t]\) 按网格点分成有限小段。在每个小段上折线斜率范数
不超过 \(M\)，故端点增量不超过 \(M\) 乘以该段长度。三角不等式求和得
\(\|y_n(t)-y_n(s)\|\le M(t-s)\)。
\end{solution}

\begin{solution}{peano-euler-residual}
在每段上 \(y_n'=f(t_k,y_n(t_k))\) 几乎处处。积分并把左端点写成
\(\kappa_n(s)\)，得到
\[
 y_n(t)=y_0+\int_{t_0}^{t}f(\kappa_n(s),y_n(\kappa_n(s)))ds.
\]
在 \(Q\) 上取距离 \(d((t,y),(s,z))=|t-s|+\|y-z\|\)，令
\(\omega_f(r)=\sup\{\|f(p)-f(q)\|:p,q\in Q,\ d(p,q)\le r\}\)。
一致连续性给 \(\omega_f(r)\to0\)（\(r\downarrow0\)）。于是
\[
 \left\|y_n(t)-y_0-\int_{t_0}^{t}f(s,y_n(s))ds\right\|
 \le h\,\omega_f((1+M)\Delta),
\]
右侧趋零。
\end{solution}

\begin{solution}{peano-two-sided-gluing}
在 \([t_0-h,t_0]\) 上以节点从 \(t_0\) 向左递推，或令
\(\widetilde y(s)=y(t_0-s)\)，把方程改写为
\(\widetilde y'=-f(t_0-s,\widetilde y)\)，应用同一右向构造。
两侧函数在 \(t_0\) 都取 \(y_0\)，故连续粘合。两侧导数趋于同一个
\(f(t_0,y_0)\)，积分方程在整个区间成立，因此粘合函数为 \(C^1\)。
\end{solution}

\begin{solution}{peano-compact-solution-set}
每个解满足 \(\|y'\|\le M\)，所以全体共同 \(M\)-Lipschitz 且取值在紧球
中。Arzelà--Ascoli 给相对紧性。若 \(y_n\) 是解且一致收敛到 \(y\)，
\(f(s,y_n(s))\to f(s,y(s))\) 一致，积分方程取极限说明 \(y\) 仍是解。
故解集闭且相对紧，即紧。
\end{solution}

\begin{solution}{peano-approximate-residual}
共同紧集条件给有界性，结合等度连续性，由 Arzelà--Ascoli 可取
\(y_{n_j}\to y\) 一致。由于 \(K\) 闭，极限图像也在 \(K\subset D\) 内。
\(f\) 在 \(K\) 上一致连续，故
\[
 \sup_s\|f(s,y_{n_j}(s))-f(s,y(s))\|\to0.
\]
在残差等式中令 \(j\to\infty\)，得到
\(y(t)=y_0+\int_{t_0}^t f(s,y(s))ds\)，所以 \(y\) 是解。
论证适用于任意一致收敛子列。若只有共同有界而没有图像紧集包含于定义域的
条件，极限可能落在域外，便不能进行上述换极限。
\end{solution}

\begin{solution}{peano-scalar-sqrt-selection}
从 \(y_n(0)=0\) 出发，每个节点的斜率 \(2\sqrt{|0|}=0\)，归纳得标准 Euler
折线恒为零。另一方面取初值 \(y_\varepsilon(0)=\varepsilon^2>0\)，对应解
\[
 y_\varepsilon(t)=\bigl(\max\{t+\varepsilon,0\}\bigr)^2.
\]
它在 \(t=-\varepsilon\) 处两侧导数都为零；其余各点直接求导，得
\(y_\varepsilon'=2\sqrt{|y_\varepsilon|}\)。故这是整条实线上的解。
在 \([-T,T]\) 上，正部函数为1-Lipschitz，因而
\[
 |y_\varepsilon(t)-(\max\{t,0\})^2|
 \le\varepsilon(2T+\varepsilon).
\]
当 \(\varepsilon\downarrow0\) 时一致趋于从零时刻立即向正向离开的解。
不能把未截断的 \((t+\varepsilon)^2\) 当作全实线上的解，它在
\(t<-\varepsilon\) 时的导数为负。不同近似方案确可选择不同解。
\end{solution}

\begin{solution}{peano-osgood-logarithm}
此时 \(\omega(s)=s^\alpha\)，而
\[
 \int_{0+}\frac{ds}{s^\alpha}
\]
在 \(\alpha\ge1\) 时发散，在 \(0<\alpha<1\) 时收敛。因此
\(\alpha\ge1\) 时全部实值零初值解唯一：正向解自动非负，反向令
\(z(s)=-y(-s)\) 再应用同一判据。\(0<\alpha<1\) 时，对每个 \(c\ge0\) 存在
\[
 y_c(t)=
 \begin{cases}
 0,&t\le c,\\
 [(1-\alpha)(t-c)]^{1/(1-\alpha)},&t>c,
 \end{cases}
\]
给出非唯一性。指数 \(1/(1-\alpha)>1\)，所以在拼接点的两侧导数均为零；
其余各点求导即得 \(y_c'=|y_c|^\alpha\)，且 \(y_c(0)=0\)。
\end{solution}

\begin{solution}{peano-continuity-not-global}
取 \(y'=y^2,y(0)=1\)。右端处处连续且局部 Lipschitz，所以局部解存在唯一；
分离变量得到 \(y(t)=1/(1-t)\)，它在 \(t\uparrow1\) 时无界，不能延拓为
经过 \(t=1\) 的有限值经典解。故连续性不保证全局存在。
\end{solution}

\begin{solution}{peano-upper-lower-solutions}
在题给的整个区间 \(I\) 上令
\[
 (Tu)(t)=y_0+\int_{t_0}^{t}f(s,u(s))ds.
\]
若 \(f\) 关于 \(y\) 单调递增，且初值满足
\(\alpha(t_0)\le y_0\le\beta(t_0)\)，则对
\(\alpha\le u\le\beta\)，
\[
 Tu-\alpha
 \ge y_0-\alpha(t_0)+\int_{t_0}^{t}
 [f(s,\alpha(s))-\alpha'(s)]ds\ge0,
\]
同理 \(Tu\le\beta\)。所以序区间不变。若从 \(\alpha,\beta\) 分别迭代，
单调性还给出两列的相应单调关系；要证明其极限和极值性质，仍需紧性或额外
连续性论证。
\end{solution}

\begin{solution}{peano-schauder-route}
Schauder 定理的一种形式是：Banach 空间中非空闭凸集 \(C\) 上的连续映射
\(T:C\to C\)，若 \(T(C)\) 相对紧，则 \(T\) 有不动点。取
\[
 C=\{u\in C(I,\R^m):\|u-y_0\|_\infty\le b\}
\]
及 Picard 积分算子，其中 \(I=[t_0-h,t_0+h]\)，\(0<h\le a\)、\(Mh\le b\)。
\(C\) 含常函数 \(y_0\)，由上确界范数不等式可见它闭且凸。
\(Mh\le b\) 给 \(T(C)\subset C\)（\(M=0\) 时自动满足）；\(f\) 在紧矩形上
一致连续，给 \(T\) 的连续性；并且 \(T(C)\) 一致有界且共同
\(M\)-Lipschitz，Arzelà--Ascoli 给相对紧性。Schauder 定理于是给不动点，
即积分方程的解。此法仍不给唯一性。
\end{solution}
